% ============================================================
\section{Fine-Tuning Protocol}
\label{sec:method}
% ============================================================

\textbf{Moirai for forecasting and anomaly scoring.}
Moirai~\cite{woo2024moirai} predicts a probabilistic distribution
over a target horizon given context.  For ETT forecasting we report
MSE on the median of 20 forecast samples at $h{=}96$ and $h{=}192$
(Appendix~\ref{app:moirai_median}).  For IoT the anomaly score is
the fraction of timesteps outside the $1{-}\alpha$ CI
($\alpha{=}0.05$) over a 32-timestep horizon.

\textbf{Fine-tuning objective.}
Condition~B uses NLL only.  Condition~C adds supervised
contrastive loss~\cite{khosla2020supcon}:
$\mathcal{L}(\theta) = \mathcal{L}_{\text{NLL}}(\theta) +
\lambda\,\mathcal{L}_{\text{SupCon}}(\mathbf{z}, \mathbf{y};\, \tau)$,
with $\lambda = 0.5$ and $\tau = 0.07$.

\textbf{Training details.}
AdamW ($\eta = 10^{-4}$, weight decay $10^{-2}$), batch size 32,
early stopping with patience 3, gradient clipping at 1.0.  ETT
fine-tuning runs 20 epochs at $n{=}500$ (10 seeds) and is swept
across $n{\in}\{500,1{\text{k}},2{\text{k}},5{\text{k}},10{\text{k}}\}$.

\textbf{uni2ts PackedStdScaler patch.\footnote{We patched an in-place
tensor mutation in \texttt{uni2ts}'s \texttt{PackedStdScaler} that
silently detaches autograd; a differentiable \texttt{torch.where}
replaces the in-place write (Appendix~\ref{app:scaler_patch}).}}
\textbf{Any Moirai fine-tuning reproduction of this paper requires the patch.}

\textbf{Representation diagnostics.}
We track representation change with two metrics:
(1)~CKA~\cite{kornblith2019cka} between pre-trained and fine-tuned
encoder representations (geometric similarity);
and (2)~weight drift ($\ell_2$ distance from pre-trained parameters).

\textbf{The value score and its baseline.}
We write $R^2_\text{task}(\text{PT}) = 1-\text{MSE}_\text{ZS}/\text{MSE}_\text{Linear}$
for the continuous pre-training value score; it is positive when the
pre-trained model beats the linear baseline, and $0.20$ marks the operating
point used as a screen rather than a boundary.
$\text{MSE}_\text{Linear}$ is one ridge map $\R^{96 \times F} \to \R^{h \times F}$
($\lambda{=}10^{-4}$) fit by least squares on the \emph{training} windows and
applied unchanged to the evaluation windows, on the train-normalised scale the
runs report, over the same window set the model is scored on (up to 300
windows). Both terms therefore see the same targets; the pre-trained model sees
$96{+}h$ steps of context against the baseline's $96$, which if anything makes
the gate a lower bound on the linear model's strength.

\textbf{Two splits, named once.} The \emph{selection} split is the set of
windows used for early stopping and checkpoint choice; the \emph{held-out} split
is disjoint from it and is where every number in this paper is scored.
$\text{B}{-}\text{D}$ is the difference in forgetting between full fine-tuning
and the frozen encoder, positive when freezing does better, and estimates the
incremental downstream value of allowing encoder-side adaptation \emph{under
this protocol}.

\textbf{Definition (degradation of a demonstrated capability).} A cell counts as
degradation when all three clauses hold: (i)~$R^2_\text{task}(\text{PT}) \ge
0.20$ on held-out windows, so there is a measured pre-trained advantage;
(ii)~full fine-tuning ends worse than not fine-tuning at all, forg.$_\text{B} >
0$, in \emph{every} seed; and (iii)~the frozen encoder ends better,
forg.$_\text{D} < 0$, in \emph{every} seed, which localises the loss to encoder
weight updates rather than to the data or the epoch budget. Clauses (ii) and
(iii) require unanimity because with three to ten seeds a majority is not
evidence of a sign. This definition is fixed for the whole paper;
\S\ref{sec:analysis:degradation} reports how many cells meet it.

\textbf{Conditions.} \textbf{B}~NLL-only fine-tune (baseline: how much does
unconstrained fine-tuning drift and forget?); \textbf{C}~NLL$+$SupCon;
\textbf{D}~frozen encoder (control: separates encoder updates from head
training and overfitting); \textbf{E}~LoRA ($r{=}8$); \textbf{F}~L2-SP and
\textbf{G}~EWC. Conditions~D and~B are the comparison every claim in this
paper turns on.

\textbf{What condition~D does and does not hold fixed.} D freezes the
encoder's own weights but leaves the input projection and mask encoding
trainable, so it is a control on \emph{encoder weight updates}, not a
guarantee of unchanged representations: its measured CKA is close to but never
exactly $1$ on Moirai ($0.969$--$0.9999$ across our cells) and falls as low as
$0.38$ on TimesFM, where the encoder stack's weight drift is exactly zero.
\textbf{H}~(strict freeze) additionally freezes those two modules and is the
only condition with CKA exactly $1.000$; we use it to check that the
$\text{B}{-}\text{D}$ comparisons do not depend on the projections being
trainable (Appendix~\ref{app:strictfreeze}).

\textbf{Dispersion.} Every ${\pm}$ in the body and in every generated table is
the standard error of the mean over seeds; two hand-assembled appendix tables
report standard deviations and say so in their captions.
